\documentclass[10pt,twoside,openright]{book}
\usepackage[UTF8,linespread=1.08]{ctex}

\usepackage[
    paperwidth=174mm,
    paperheight=247mm,
    textwidth=29pc,
    top=25mm,
    bottom=25mm,
    centering,
]{geometry}

% Page headers with rule (matching EngC: rule extends 4pc into left margin)
\usepackage{fancyhdr}
\pagestyle{fancy}
\fancyhfoffset[L]{4pc}
\fancyhf{}
\fancyhead[LE]{\thepage}
\fancyhead[RE]{\nouppercase{\leftmark}}
\fancyhead[LO]{\nouppercase{\rightmark}}
\fancyhead[RO]{\thepage}
\renewcommand{\headrulewidth}{0.4pt}

% Font packages
\usepackage[utf8]{inputenc}
\usepackage{bbold}
\usepackage{pifont}

% Math packages
\usepackage{amsmath}
\usepackage{amssymb}

% Figures and colors
\usepackage{xcolor}
\usepackage{graphicx}
\usepackage{caption}
\usepackage{float}
\usepackage{subcaption}

% Diagrams and images
\usepackage{tikz} % General diagram software
\usepackage{tikz-cd} % Commutative diagrams
\usepackage{pgfplots} % PGF plots; used to plot functions
\usetikzlibrary{automata, arrows, angles, calc, intersections, math,
decorations.pathmorphing, decorations.markings}
\usetikzlibrary{positioning}
%% the following commands are needed for some matlab2tikz features
\usetikzlibrary{plotmarks}
\usetikzlibrary{arrows.meta}
\pgfplotsset{compat=1.16}
\usepgfplotslibrary{groupplots} % This lets you put multiple
% functions together in a single plot
\usepgfplotslibrary{patchplots}

% Tables
\usepackage{booktabs}
\usepackage{colortbl}
\usepackage{multirow}

% References and bibliography
\usepackage[
    backend=biber,
    style=alphabetic,
    sorting=nyt,
    doi=false,
    hyperref=true,
    isbn=false,
    sortcites,
    url=false,
    maxbibnames=99,
    maxalphanames=3,
    minalphanames=3,
    maxcitenames=2,
    mincitenames=1,
    backref=false,
    natbib=true,
]{biblatex}
\renewbibmacro{in:}{}
\AtEveryBibitem{%
    \clearlist{language}%
}
% Hack to accommodate texlab users
\providecommand{\toplevelprefix}{.}  % necessary for subfile
% compilation to work with bibliography
\makeatletter
\@ifundefined{thiscommandisneverdefined}{
    \addbibresource{\toplevelprefix/book-references.bib}
}{
    \addbibresource{book-references.bib}
}
\makeatother

\usepackage{hyperref}
\hypersetup{
    colorlinks=true,
    citecolor={blue!80!black},
}
\usepackage{cleveref}
\crefname{assumption}{assumption}{assumptions}
\Crefname{assumption}{Assumption}{Assumptions}
\crefname{exercise}{exercise}{exercises}
\Crefname{exercise}{Exercise}{Exercises}

\hyphenation{line-break line-breaks docu-ment triangle cambridge amsthdoc
cambridgemods baseline-skip author authors cambridgestyle en-vir-on-ment polar}

\setcounter{tocdepth}{2}% list subsections in the table of contents
% Algorithm packages -- should go after hyperref
\usepackage{algorithm}
% \usepackage[spaceRequire=false]{algpseudocodex}
\usepackage{algpseudocode}

% Math formatting
\usepackage{mathtools}
\usepackage{grffile}

% Text formatting
\usepackage[shortlabels]{enumitem}
\setlist[enumerate]{label=\arabic*., leftmargin=*}
\usepackage{fancyvrb}
\usepackage{soul}
\VerbatimFootnotes
\allowdisplaybreaks

% Custom packages for the book
\usepackage{math-macros}  % Custom package
\usepackage{math-theorems}  % Custom package
\usepackage{book-macros}  % Custom package for macros in the book

% Add subfiles --- each chapter can now be compiled by itself!
\usepackage{subfiles}

% Dotted leaders for chapter TOC entries (enabled after frontmatter)
\makeatletter
\newif\if@chapterdots \@chapterdotsfalse
\renewcommand*\l@chapter[2]{%
  \ifnum \c@tocdepth >\m@ne
    \addpenalty{-\@highpenalty}%
    \vskip 1.0em \@plus\p@
    \setlength\@tempdima{1.5em}%
    \begingroup
      \parindent \z@ \rightskip \@pnumwidth
      \parfillskip -\@pnumwidth
      \leavevmode \bfseries
      \advance\leftskip\@tempdima
      \hskip -\leftskip
      #1\nobreak
      \if@chapterdots
        \leaders\hbox{$\m@th\mkern \@dotsep mu\hbox{.}\mkern \@dotsep mu$}\hfill
      \else
        \hfil
      \fi
      \nobreak\hbox to\@pnumwidth{\hss #2}\par
      \penalty\@highpenalty
    \endgroup
  \fi}
\makeatother

% Title page customization using titling package
\usepackage{titling}
\pretitle{\begin{center}\LARGE\bfseries}
\posttitle{\end{center}\vskip 0.5em}
\preauthor{\begin{center}\large}
\postauthor{\end{center}}
\predate{\begin{center}\large}
\postdate{\end{center}}

% Customize vertical spacing
\setlength{\droptitle}{-5em}

% Add footnote content below date
\renewcommand{\maketitlehookd}{%
  \vspace{4em}\footnotesize $^*$\,These authors contributed equally.
}

\begin{document}

 
\title{深度表征学习的原理与实践\\
\centerline{\LARGE 结构化记忆的数学理论}
%\vspace{0.1in}\\ 
%\makebox[\textwidth][c]{Structured Memory: Principles and Practice of Deep Representation Learning}
}

\author{\vspace{0.5cm}
\textbf{Sam Buchanan}$^*$, 
        University of California, Berkeley \& TTIC  \\[\baselineskip]
        \textbf{Druv Pai}$^*$, University of California, Berkeley
        \\[\baselineskip]
        \textbf{Peng Wang}, 
        University of Macau \& University of Michigan  \\[\baselineskip]
        \textbf{Yi Ma}, 
        University of Hong Kong \& University of California, Berkeley
	\vspace{5cm}\\
	\footnotesize
	This manuscript is prepared for a course on the principles of deep learning
  and intelligence in general, to be taught by the authors at the respective
  institutions. This pre-publication version is free to view and download for
  personal use only, and is not for redistribution, re-sale or use in derivative
  works.\\  Copyright \Pisymbol{psy}{227} 2025 Sam Buchanan, Druv Pai, Peng Wang, Yi Ma
	%\footnotesize Copyright \Pisymbol{psy}{227}2020 Reserved \\ No
	%parts of this draft may be reproduced or distributed without written permission
	%from the authors.\\
}

\date{\large 2026年3月1日}
% \date{\large \today}

\frontmatter
\titlepage
\thispagestyle{empty}
\maketitle

\addcontentsline{toc}{chapter}{序言}
\subfile{chapters/preface/preface_zh.tex}

\addcontentsline{toc}{chapter}{\numberline{}第二版序言}
\subfile{chapters/preface/preface-v2_zh.tex}

\addcontentsline{toc}{chapter}{\numberline{}开源声明}
\subfile{chapters/declaration/declaration_zh.tex}

\addcontentsline{toc}{chapter}{\numberline{}致谢}
\subfile{chapters/acknowledgment/acknowledgment_zh.tex}

\addcontentsline{toc}{chapter}{\numberline{}符号表示}
\subfile{chapters/notation_zh.tex}

\addtocontents{toc}{\protect\@chapterdotstrue}

\tableofcontents

\setcounter{page}{1}
\pagenumbering{arabic}

\mainmatter

\subfile{chapters/chapter1/introduction_zh}

\subfile{chapters/chapter2/classic-models_zh}

\subfile{chapters/chapter3/denoising_zh}

\subfile{chapters/chapter4/lossy-compression_zh}

\subfile{chapters/chapter5/deep-networks_zh}

\subfile{chapters/chapter6/consistent-representations_zh}

\subfile{chapters/chapter7/conditional-inference_zh}

\subfile{chapters/chapter8/applications_zh}

\subfile{chapters/chapter9/future_zh}

\newpage\appendix
\part*{Appendices}
\addcontentsline{toc}{part}{Appendices}

\subfile{chapters/appendixA/optimization_zh}

\subfile{chapters/appendixB/entropy_zh}

\addtocontents{toc}{\vspace{\baselineskip}}

%\addcontentsline{toc}{chapter}{Bibliography}
\printbibliography

\end{document}
